\documentclass{article}
\usepackage[T1]{fontenc}
\usepackage{lmodern}
\usepackage{graphicx}
\usepackage{amsmath,amssymb}
\usepackage[a4paper,margin=2cm]{geometry}
\usepackage[absolute,overlay]{textpos}
\setlength{\TPHorizModule}{1mm}
\setlength{\TPVertModule}{1mm}
\usepackage[indonesian]{babel}
\usepackage{hyperref}
\setlength{\emergencystretch}{2em}
% unit: Halaman 13

\begin{document}

% === Halaman 13 (llm) ===

\begin{itemize}
    \item Kunci untuk OTP harus seluruhnya acak sejati dan sepanjang pesan.
    \item Bagaimana jika kunci diambil dari teks yang panjang (misalnya tulisan di dalam novel, buku, berita, dan sebagainya)?
    \begin{itemize}
        \item ini bukan lagi OTP (sebab tulisan di buku/novel/berita bukan acak)
        \item tidak menghasilkan \textit{perfect secrecy}
        \item dapat dipecahkan
    \end{itemize}
    \item Kunci di dalam OTP hanya dipakai sekali dan tidak pernah digunakan kembali. Bagaimana jika kunci dipakai untuk kedua kalinya?
    \begin{itemize}
        \item ia bukan lagi \textit{one-time pad}, tetapi \textit{two-time pad}
        \item tidak aman
    \end{itemize}
\end{itemize}

\end{document}
